\documentclass[11pt,a4paper]{article}

% ============================================================
% DelegationChain paper — preamble
% Packages, settings, custom commands.
% ============================================================

% --- Document basics ---
\usepackage[a4paper, margin=2.5cm]{geometry}
\usepackage[T1]{fontenc}
\usepackage[utf8]{inputenc}
\usepackage{lmodern}
\usepackage{microtype}

% --- Math ---
\usepackage{amsmath, amsthm, amssymb}
\usepackage{mathtools}

% --- Tables and figures ---
\usepackage{booktabs}
\usepackage{array}
\usepackage{graphicx}
\usepackage{caption}
\captionsetup{font=small, labelfont=bf}

% --- Code listings ---
\usepackage{listings}
\usepackage{xcolor}
\definecolor{codebg}{HTML}{F7F7F7}
\definecolor{codeborder}{HTML}{D0D0D0}
\lstset{
    basicstyle=\ttfamily\small,
    backgroundcolor=\color{codebg},
    frame=single,
    rulecolor=\color{codeborder},
    breaklines=true,
    keepspaces=true,
    columns=fullflexible,
    xleftmargin=4pt,
    xrightmargin=4pt,
    aboveskip=8pt,
    belowskip=8pt,
}

% --- Cross-references ---
\usepackage{hyperref}
\hypersetup{
    colorlinks=true,
    linkcolor=black,
    urlcolor=blue!50!black,
    citecolor=blue!50!black,
    pdftitle={DelegationChain},
    pdfauthor={Daniyar Amangeldin},
}
\usepackage[capitalize]{cleveref}

% --- Bibliography ---
\usepackage[
    backend=biber,
    style=authoryear,
    sorting=nyt,
    natbib=true,
    maxcitenames=2,
    mincitenames=1,
    maxbibnames=99,
]{biblatex}
\addbibresource{references.bib}

% --- Theorem environments ---
\newtheorem{theorem}{Theorem}
\newtheorem{lemma}[theorem]{Lemma}
\newtheorem{proposition}{Proposition}
\theoremstyle{definition}
\newtheorem{definition}{Definition}
\theoremstyle{remark}
\newtheorem*{remark}{Remark}

% --- Spacing ---
\setlength{\parskip}{0.4em}
\setlength{\parindent}{1.2em}

% --- Custom commands ---
\newcommand{\G}[1]{\mathbb{G}_{#1}}
\newcommand{\Gt}{\mathbb{G}_T}
\newcommand{\Zq}{\mathbb{Z}_q}
\newcommand{\pair}[2]{e(#1, #2)}
\newcommand{\Sign}{\mathsf{Sign}}
\newcommand{\Verify}{\mathsf{Verify}}
\newcommand{\HashG}{\mathsf{HashToG}_2}
\newcommand{\Canon}{\mathsf{Canon}}
\newcommand{\sk}{\mathit{sk}}
\newcommand{\pk}{\mathit{pk}}

% Threat label macro
\newcommand{\threat}[2]{\textbf{T#1:} \emph{#2}}

\title{DelegationChain: Aggregatable Capability Chains for \\
       Cross-Organizational Agent Authorization}

\author{
    Daniyar Nurmukhamet\thanks{Contact: \texttt{nurmuhamet.d@edu.ektu.kz}. Working draft;
    not for citation.} \\
    School of Digital Technologies and Artificial Intelligence \\
    D. Serikbayev East Kazakhstan Technical University \\
    \"Oskemen, Kazakhstan
}

\date{\today}

\begin{document}

\maketitle

\begin{abstract}
Autonomous software agents built on large language models increasingly invoke
tools and services across organizational boundaries, frequently via the Model
Context Protocol (MCP). When an agent at one organization invokes a tool at
another, the receiving organization currently has no cryptographically
verifiable guarantee that the invocation was authorized, that the specific
parameter values were sanctioned rather than just the action type, or that the
authorization chain leading to the call was not tampered with. Existing
authorization protocols address subsets of these properties but do not
jointly satisfy them. We propose DelegationChain, a protocol providing
parameter-bound, aggregatable, offline-verifiable authorization proofs for
cross-organizational agent invocations. The protocol uses BLS12-381 signatures
with proof-of-possession, a recursive chain digest construction enabling
compact aggregate signatures across $N$-hop delegations, a typed declarative
policy expression language with decidable subset relations, and a federated
identity registry model that permits interoperation between mutually
distrusting organizations without a shared trust anchor. We define a threat
model covering seventeen attack classes, describe the protocol's
cryptographic and operational components, and present a preliminary security
analysis. We discuss limitations, including the absence at this stage of an
empirical performance evaluation and machine-checked security proofs, and
outline the path to a full evaluation including symbolic protocol verification
and a Rust reference implementation.
\end{abstract}

% ============================================================
\section{Introduction}
\label{sec:introduction}

Software agents built on large language models are increasingly deployed in
configurations where they invoke external tools and delegate sub-tasks to
other agents. The Model Context Protocol \citep{anthropic-mcp}, introduced in
late 2024, has emerged as a widely-adopted standard for this interaction
pattern, with reported adoption across thousands of integrating services.
As deployment scales, an architectural pattern is becoming common: agents
operated by one organization invoke tools or delegate work to agents
operated by another. Industry observers have called the resulting
multi-organizational ecosystem the inter-agent economy.

Cross-organizational agent invocation creates an authorization problem that
existing infrastructure does not adequately address. Consider an agent at
Organization A, acting on behalf of a human principal, that needs to invoke
a payment-processing tool at Organization B. Organization B's tool must
decide whether to honor the request. To do so safely, B needs verifiable
answers to four questions: \emph{(i)} did the request actually originate
from an authorized agent at A, rather than from an attacker impersonating
that agent; \emph{(ii)} was the agent operating within the scope its
human principals authorized, rather than acting beyond its delegated
authority; \emph{(iii)} were the specific parameter values supplied --- the
amount, the recipient, the file path --- sanctioned by the authorization
chain, or has an intermediate party substituted them; and \emph{(iv)} has
the authorization chain itself been preserved intact, or has some step in
the chain been forged, dropped, or reordered.

Existing protocols address subsets of these questions but do not jointly
satisfy them. OAuth 2.1 \citep{oauth21} proves an agent's identity to a
resource server but does not bind authorization to specific parameter
values, requires a synchronous call to the authorization server on the
verification path, and cannot represent attenuated delegated authority across
organizational boundaries. Macaroons \citep{birgisson2014macaroons} provide
attenuated capability delegation with policy caveats, but use HMAC-based
verification with shared secrets, which does not extend cleanly across
organizations that do not share infrastructure. SPIFFE \citep{spiffe} provides
workload identity within a trust domain but does not include parameter binding
or delegation chains. UCAN \citep{ucan-spec} extends capability delegation to
public-key settings but is designed for Web3 contexts and does not address
MCP integration, BLS aggregation, or policy language formalization. Recent
MCP-specific proposals such as AGENTBOUND \citep{buhler2025agentbound} and
the access-control mechanisms identified by \citet{huang2026inch} and
\citet{maloyan2026breaking} provide single-organization runtime enforcement
or measure the gap, but do not provide cross-organizational cryptographic
proofs. The Agent Identity Protocol \citep{prakash2026aip} introduces
invocation-bound capability tokens for cross-protocol delegation but does not
incorporate aggregatable signatures, the formal policy language we present,
or the federated identity infrastructure.

This paper proposes DelegationChain, an authorization protocol designed
specifically for cross-organizational agent invocation. We describe a
protocol design intended to combine, in a single coherent system, several
properties that prior systems have addressed individually: public-key
cryptographic verification across organizations without a shared root of
trust; binding of authorization to exact parameter values, not merely action
types; compact aggregate signatures across multi-hop delegation chains; and
verification that requires no synchronous calls to the issuing authority
during the dominant path. The protocol is constructed on BLS12-381
signatures \citep{boneh2001bls, ietf-bls-signature} with proof-of-possession
registration, a recursive chain digest construction enabling chain integrity
under aggregation, a typed declarative policy language with decidability
guarantees, and a federated organizational registry model.

\paragraph{Contributions.}
This paper makes four contributions. First, we define a threat model for
cross-organizational agent invocation that explicitly covers seventeen
distinct attack classes spanning protocol-level attacks, identity
infrastructure attacks, and threats specific to LLM-based agent systems
such as prompt injection within permitted policy
(\cref{sec:threat-model}). Second, we specify the protocol's cryptographic
core, wire format, verifier algorithm, and identity infrastructure
(\cref{sec:protocol,sec:identity}) to a level of detail sufficient for an
independent reference implementation. Third, we introduce a typed
declarative policy expression language with a formal grammar, decidable
subset relation, and proof of evaluator termination (\cref{sec:policy}),
which we argue is sufficiently expressive for realistic enterprise
authorization while remaining free of the injection and denial-of-service
hazards of general-purpose policy languages. Fourth, we provide a
preliminary security analysis (\cref{sec:security}) reducing the protocol's
core security properties to standard cryptographic assumptions.

\paragraph{Scope and limitations.}
This paper is a design contribution. We do not yet present empirical
performance measurements from a reference implementation, machine-checked
security proofs, or results from symbolic protocol verification tools such
as Tamarin or ProVerif. Each of these is in progress and is required before
the protocol is ready for production use. We discuss the planned evaluation in
\cref{sec:limitations}. The contribution of this paper is the design and
its preliminary analysis; the full evaluation is future work.

\paragraph{Paper organization.}
\Cref{sec:background} reviews necessary background on BLS signatures and
the MCP architecture. \Cref{sec:threat-model} defines the threat model.
\Cref{sec:protocol} describes the cryptographic core: chain digest
construction, encoding, approval receipts, and the verifier algorithm.
\Cref{sec:identity} describes the federated identity infrastructure.
\Cref{sec:policy} introduces the policy expression language.
\Cref{sec:security} sketches the security argument. \Cref{sec:discussion}
discusses related work and limitations. \Cref{sec:conclusion} concludes.

% ============================================================
\section{Background}
\label{sec:background}

This section reviews the cryptographic and protocol background required to
follow the design. Readers familiar with pairing-based cryptography and MCP
may skim.

\subsection{Model Context Protocol}
\label{sec:bg-mcp}

The Model Context Protocol (MCP) \citep{anthropic-mcp} is a JSON-RPC-based
protocol standardizing how AI agents discover and invoke external tools.
An MCP \emph{host} (typically an LLM-based agent) communicates with one or
more MCP \emph{servers}, each exposing a set of tools. Tools are described
by name, a typed parameter schema, and documentation. The host requests a
tool invocation by sending a JSON-RPC message specifying the tool name and
parameter values; the server executes the tool and returns the result.

MCP's design assumes that the host and server are mutually trusted, an
assumption appropriate for the single-organization deployments common in
its early adoption. As MCP deployments scale to cross-organizational
configurations, this assumption becomes untenable. A server invoked by an
unknown agent has no cryptographic basis on which to decide whether the
invocation is authorized. The protocol's current security guidance relies
on transport-layer authentication (typically OAuth) which addresses caller
identity but not the broader authorization questions enumerated in
\cref{sec:introduction}. Empirical measurements of deployed MCP servers
\citep{huang2026inch} have found a substantial fraction lacking even basic
caller identity verification.

\subsection{BLS Signatures and Aggregation}
\label{sec:bg-bls}

The Boneh-Lynn-Shacham (BLS) signature scheme \citep{boneh2001bls} is a
signature scheme constructed from bilinear pairings on elliptic curves.
On the BLS12-381 curve \citep{ietf-bls-signature}, signatures consist of
a single point in $\G{2}$ (96 bytes compressed) and public keys consist
of a single point in $\G{1}$ (48 bytes compressed). Signing requires one
scalar multiplication; verification requires evaluation of a bilinear
pairing.

The property that makes BLS signatures particularly attractive for our
setting is \emph{signature aggregation}. Given signatures
$\sigma_1, \ldots, \sigma_n$ produced by distinct signers on distinct
messages, their sum $\sigma_{\text{agg}} = \sum_i \sigma_i$ is itself a
valid object in $\G{2}$, and its validity with respect to the signing
public keys and messages can be checked via a single multi-pairing
operation \citep{boneh2018compact}. This aggregation has two consequences
relevant to our design.

First, the bandwidth required to transmit $n$ signatures is reduced to a
single $\G{2}$ element, regardless of $n$. For our delegation chains, this
means a chain of arbitrary length carries a fixed 96-byte aggregate
signature rather than $96n$ bytes of individual signatures.

Second, the computational cost of verification is reduced relative to
verifying $n$ signatures independently. We address the precise cost in
\cref{sec:protocol-verification}; here we note that the cost is not, as is
sometimes stated, asymptotically $O(1)$ in the number of signers when the
signers sign distinct messages. The honest characterization, which we
adopt throughout, is that aggregation provides significant constant-factor
savings via Miller loop sharing in the underlying multi-pairing
implementation.

BLS aggregation is vulnerable to rogue-key attacks unless the registration
of public keys requires proof-of-possession of the corresponding secret
keys \citep{boneh2018compact}. We address this in \cref{sec:identity-pop}.

\subsection{Capability-Based Authorization}
\label{sec:bg-capabilities}

The protocol design draws on the capability-based security tradition
\citep{miller2006robust}. In a capability system, authority is represented
by unforgeable tokens (capabilities) that the holder can present to obtain
access. Capabilities can be \emph{attenuated}: a holder may produce a
restricted version of a capability and delegate the restricted version to
another party, without being able to grant authority exceeding what they
themselves hold.

Macaroons \citep{birgisson2014macaroons} apply this model to web
authorization. A macaroon is an HMAC chain in which each link adds
``caveats'' --- predicates restricting the macaroon's use. Verification
recomputes the HMAC chain using a shared secret known to the verifier.
This works well within a single trust domain but does not extend cleanly
across organizational boundaries, because organizations do not generally
share HMAC secrets.

UCAN \citep{ucan-spec} extends capability attenuation to public-key
settings using ECDSA or Ed25519 signatures and a JWT-like token format.
UCAN is designed for Web3 contexts and lacks features specific to our
setting, including signature aggregation, formal policy language
specification, and MCP transport binding.

DelegationChain inherits the conceptual framework of attenuated
capabilities from this tradition. Our contribution is the combination of
public-key verification (extending Macaroons' attenuation model across
organizational boundaries), signature aggregation (improving over the
per-hop linear growth of UCAN tokens), an explicit policy language with
decidability guarantees, and a federated identity infrastructure.

% ============================================================
\section{Threat Model}
\label{sec:threat-model}

This section defines the threat model under which DelegationChain operates.
We state the protocol's assets, the assumed adversary capabilities, the
threats we defend against, the threats we explicitly do not defend against,
and the security properties we seek to provide.

\subsection{Assets and Principals}
\label{sec:threat-assets}

The assets we seek to protect are: \emph{(i)} the integrity of tool
invocations, in the sense that an accepted invocation faithfully
represents the authorization granted by the originating principals;
\emph{(ii)} the provenance of authorization chains, in the sense that the
chain delivered to a verifier accurately records the parties that authorized
each step; and \emph{(iii)} the integrity of audit records derived from
accepted chains.

The principals participating in the protocol are: the \emph{human operator},
who initiates a request and whose authority is the ultimate source of
permissions; the \emph{orchestrator agent}, which receives the initial
authorization and may delegate to other agents; \emph{specialist agents},
which receive delegated authority and either invoke tools directly or
delegate further; the \emph{tool} or \emph{MCP server}, which receives
authorization chains and decides whether to honor invocations; the
\emph{signing service}, a hosted component that holds an agent's private
key and applies policy enforcement before signing; the \emph{approval
service}, which obtains human approval for actions requiring it and signs
approval receipts; and the \emph{registry}, which issues identity
certificates binding agents and services to public keys.

\subsection{Adversary Capabilities}
\label{sec:threat-adversary}

We model a probabilistic polynomial-time adversary with the following
capabilities. The adversary can observe and modify all network traffic
between principals (a standard Dolev-Yao adversary on the network). The
adversary can compromise an agent's runtime environment, gaining the
ability to instruct the signing service to sign messages of the
adversary's choice within the scope permitted by policy. The adversary can
inject content into agent inputs to attempt prompt injection attacks. The
adversary cannot, by default, compromise hardware security modules, the
signing service infrastructure, the approval service infrastructure, the
registry, or human approvers' devices; threats involving these
compromises are addressed separately and explicitly.

\subsection{Threats Within Scope}
\label{sec:threat-within-scope}

The protocol is designed to defend against the following threats. Threats
are numbered for cross-reference in the security analysis
(\cref{sec:security}).

\threat{1}{Token forgery.} An adversary attempts to produce a valid
authorization token without holding the corresponding signing key.

\threat{2}{Parameter substitution.} An adversary, positioned between a
signer and a verifier, attempts to alter parameter values in an
authorized invocation (e.g., changing the destination of a transfer)
while preserving the appearance of valid authorization.

\threat{3}{Scope escalation.} An agent attempts to delegate more authority
than it was itself granted.

\threat{4}{Chain forgery and reordering.} An adversary attempts to
construct an authorization chain that was not produced by the legitimate
delegation sequence, including reordering hops, dropping hops, or
inserting unauthorized hops.

\threat{5}{Prompt injection within policy.} An adversary, controlling
content fed to an agent, induces the agent to take actions that are
within the agent's permitted policy but contrary to the principal's
intent.

\threat{6}{Replay across sessions.} An adversary captures a valid
authorization chain and attempts to use it after its intended validity
window has expired.

\threat{7}{Replay within validity window.} An adversary captures a valid
authorization chain and attempts to use it more than once within its
validity window.

\threat{8}{Confused deputy in delegation.} An agent is induced to sign a
delegation token favoring an adversary by manipulation of the inputs the
agent uses to construct delegation decisions.

\threat{9}{Cross-context token reuse.} An adversary attempts to reuse a
token issued in one context (e.g., for a delegation) as if it had been
issued in another context (e.g., for an invocation).

\threat{10}{Compromised approval service.} An adversary obtains the
signing key of an approval service and issues fraudulent approvals.

\threat{11}{Approval routing attack.} An agent routes approval requests
to a permissive approval service when policy requires a stricter one.

\threat{12}{Approval reuse across invocations.} An adversary attempts to
reuse a valid approval receipt with an invocation different from the one
approved.

\threat{13}{Stale approval after parameter modification.} An adversary
attempts to modify invocation parameters after approval has been
obtained.

\threat{14}{Rogue-key registration.} An adversary attempts to register a
public key for which it does not hold the secret key, in order to inject
controlled keys into aggregate signatures.

\threat{15}{Compromised registry signing key.} An adversary obtains a
registry's root signing key and issues fraudulent certificates.

\threat{16}{Revocation suppression.} An adversary suppresses revocation
information from reaching verifiers, extending the window during which
compromised credentials remain accepted.

\threat{17}{Certificate substitution.} An adversary substitutes a
different certificate for the same identity during the registration or
resolution process.

\subsection{Threats Outside Scope}
\label{sec:threat-outside}

The following threats are explicitly outside our scope. Defending against
them requires mechanisms orthogonal to the protocol.

A fully compromised signing service can produce signatures on any
messages within an agent's policy; the protocol cannot prevent this and
relies on operational hardening (HSM-based key custody, monitored
access). A fully compromised approval service can produce fraudulent
approvals; mitigation is via transparency log monitoring (\cref{sec:identity})
and rotation. Side-channel attacks on cryptographic implementations are
out of scope; we rely on the security properties of well-audited BLS
libraries such as \texttt{blst}. Attacks on the human approver's device or
authentication mechanism are out of scope; the protocol assumes a trusted
path between the approval service and the human. Denial-of-service attacks
on network infrastructure are out of scope. Post-quantum security is out
of scope for this version; BLS12-381 is not post-quantum secure, and
migration to a post-quantum scheme is left to future work.

\subsection{Security Properties Sought}
\label{sec:threat-properties}

We seek to provide the following security properties, each of which is
formally stated and argued in \cref{sec:security}:

\begin{description}
    \item[Authenticity.] Authorization tokens are unforgeable: no PPT
    adversary lacking access to a signing key can produce a token that
    verifies under the corresponding public key, except with negligible
    probability.
    \item[Parameter binding.] A token authorizing an invocation with
    parameter values $P$ is invalid for any parameter values $P' \neq P$.
    \item[Chain integrity.] The aggregate signature over a delegation
    chain covers the entire chain content; modification, reordering,
    truncation, or extension of the chain invalidates the aggregate.
    \item[Scope monotonicity.] In any valid chain, the scope at hop
    $k+1$ is a subset of the scope at hop $k$.
    \item[Non-repudiation.] Neither party to a verified action can
    plausibly deny it, assuming honest signing key management.
    \item[Bounded revocation latency.] When a key compromise is reported
    to the registry, the maximum window during which adversary-produced
    tokens remain accepted is bounded by a configurable parameter
    (default: 120 seconds).
\end{description}

% ============================================================
\section{Protocol Design}
\label{sec:protocol}

This section presents the cryptographic core of DelegationChain: the
participants, the token types, the chain digest construction enabling
aggregate signatures across delegations, the encoding of parameter-bound
authorization, the handling of approval receipts, and the verifier
algorithm. We describe the design at a level of detail sufficient to
support the security analysis of \cref{sec:security} and to motivate the
implementation effort; a separate specification document (referenced in
\cref{sec:discussion}) provides the bit-level details required for
interoperable implementations.

\subsection{Participants and Tokens}
\label{sec:protocol-participants}

A DelegationChain authorization chain has the following participants, in
the order they contribute to the chain. The \emph{issuer}~$I$ is the
registry or identity authority that issues an initial session token to
the orchestrator agent at the start of an authorized session. The
\emph{orchestrator}~$A_1$ is the agent that receives the session token
and initiates the workflow. Intermediate \emph{specialist
agents}~$A_2, \ldots, A_{N-1}$ may be delegated sub-tasks. The
\emph{invoking agent}~$A_N$ is the agent that finally invokes a tool at
the verifier. The \emph{verifier}~$V$ is the receiving tool or MCP
server. When policy requires human approval, an \emph{approval service}
contributes a separately-signed receipt.

The chain contains three token types. The \emph{session token},
signed by the issuer, authorizes the orchestrator to operate within a
scope for a bounded duration. \emph{Delegation tokens}, signed by
intermediate agents $A_k$ for $1 \le k < N$, attenuate authority and
pass it to the next agent. The \emph{invocation token}, signed by the
final agent $A_N$, binds the invocation to specific tool, action, and
parameter values. The chain delivered to the verifier consists of all
token bodies plus a single aggregate signature in $\G{2}$.

\subsection{Cryptographic Primitives}
\label{sec:protocol-primitives}

We use BLS12-381 with the minimal-pubkey-size variant: public keys are
points in $\G{1}$ (48 bytes compressed), signatures are points in $\G{2}$
(96 bytes compressed), and the pairing maps $\G{1} \times \G{2}$ to
$\Gt$. Hash-to-curve uses RFC 9380 with suite
\texttt{BLS12381G2\_XMD:SHA-256\_SSWU\_RO\_}. SHA-256 is used for all
non-curve hashing, including the chain digest construction.

Signing key $\sk \in \Zq$, public key $\pk = \sk \cdot g_1 \in \G{1}$.
For message $m \in \{0,1\}^*$, $\Sign(\sk, m) = \sk \cdot \HashG(m) \in
\G{2}$. Verification checks $\pair{g_1}{\sigma} \overset{?}{=}
\pair{\pk}{\HashG(m)}$. For signatures
$\sigma_1, \ldots, \sigma_n$ from distinct signers on distinct messages,
the aggregate $\sigma_{\text{agg}} = \sum_i \sigma_i \in \G{2}$ is
verified by a multi-pairing check (\cref{sec:protocol-verification}).

All participants register public keys through a process that includes
proof-of-possession (\cref{sec:identity-pop}), preventing rogue-key
attacks against the aggregate.

\subsection{Chain Digest Construction}
\label{sec:protocol-digest}

The central cryptographic construction is a recursive digest chain
binding each hop to all prior hops. We use domain separation tags to
ensure tokens of different types cannot be confused (preventing
cross-context reuse, threat T9). Tags are 8-byte ASCII strings with a
null terminator: \texttt{TAG\_SES} for sessions, \texttt{TAG\_DEL} for
delegations, \texttt{TAG\_INV} for invocations.

Let $\Canon(\cdot)$ denote canonical encoding (deterministic CBOR per
RFC 8949 §4.2.1, with the additional rules summarized in
\cref{sec:protocol-encoding}).

\paragraph{Session message.}
The issuer signs:
\begin{equation}
m_0 = H(\text{TAG\_SES} \,\|\, \Canon(\text{SessionBody}))
\end{equation}
where SessionBody contains the orchestrator's identifier and public key,
a session identifier, a hash committing to the governing policy, a
scope expression, issue and expiry timestamps, the issuer's identifier,
and a fresh nonce.

\paragraph{Delegation messages.}
For each intermediate hop $A_k$ with $1 \le k < N$:
\begin{equation}
m_k = H(\text{TAG\_DEL} \,\|\, m_{k-1} \,\|\, \Canon(\text{DelegationBody}_k))
\end{equation}
where DelegationBody$_k$ contains the delegator's public key, the
delegatee's identifier and public key, the attenuated sub-scope, a hop
index, the parent session identifier, expiry, and a fresh nonce. The
critical property is that $m_k$ commits to $m_{k-1}$ via the hash,
recursively committing to the entire chain history.

\paragraph{Invocation message.}
The final hop $A_N$ signs:
\begin{equation}
m_N = H(\text{TAG\_INV} \,\|\, m_{N-1} \,\|\, \Canon(\text{InvocationBody}))
\end{equation}
where InvocationBody contains the invoker's public key, the target tool
and action identifiers, the canonically encoded parameter values, a
hash of the parameters (defense-in-depth against encoding bugs), an
optional approval receipt, and timing and nonce fields.

\paragraph{Aggregate signature.}
Each participant signs its corresponding message with its private key,
producing $\sigma_k = \Sign(\sk_k, m_k)$ for $k = 0, 1, \ldots, N$. The
aggregate signature transmitted to the verifier is
\begin{equation}
\sigma_{\text{agg}} = \sigma_0 + \sigma_1 + \cdots + \sigma_N \in \G{2},
\end{equation}
which has constant size 96 bytes regardless of $N$.

\paragraph{Role distinguishability by position.}
The verifier reconstructs the $m_k$ values from received token bodies
using the tag determined by token position: \texttt{TAG\_SES} for the
session token, \texttt{TAG\_DEL} for each intermediate delegation
token, and \texttt{TAG\_INV} for the invocation token. The verifier
\emph{rejects} chains where the position-derived tag does not match
the token kind --- preventing chain-extension attacks that would attempt
to reuse an invocation token as a delegation, or vice versa.

\subsection{Parameter Binding}
\label{sec:protocol-params}

Parameter binding is achieved by including the canonical encoding of
parameter values directly in the message $m_N$ that the invoking agent
signs. Any modification to parameter values --- changing an amount,
substituting a recipient, altering a file path --- changes the
canonical encoding, which changes $m_N$, which invalidates the
signature contribution $\sigma_N$ to the aggregate.

To support nested parameter structures (e.g.,
\texttt{order.payment.amount}), parameter values are encoded as a
nested CBOR map with text-string keys sorted lexicographically by
encoded form (per RFC 8949 §4.2.1). Parameter values are restricted to
unsigned and signed integers, byte strings, NFC-normalized text
strings, booleans, null, and nested maps and arrays of the same. Float
values and tagged values are disallowed in parameter encodings, to
eliminate edge cases around floating-point comparison and the
introduction of unexpected types.

The InvocationBody also includes \texttt{params\_hash}, a SHA-256
digest of the canonical parameter encoding. The verifier recomputes
this hash from the encoded parameters and rejects any chain where it
does not match. This is redundant with the signature check but serves
as defense-in-depth against subtle encoding bugs in implementations.

\subsection{Approval Receipts}
\label{sec:protocol-approvals}

When policy requires human approval for an action, the InvocationBody
carries an \emph{approval receipt} signed by a registered approval
service. The receipt binds the human's approval decision to the
specific invocation being approved.

The flow is as follows. The agent constructs the InvocationBody
\emph{except} for the approval receipt field, computes a digest over
this preliminary body, and submits the digest along with the full body
to the required approval service. The service presents the proposed
invocation to a human approver. On approval, the service signs an
\emph{approval body} containing the service's identifier and public
key, the invocation digest, opaque attestation data identifying the
approver, and timing fields. The signed receipt is returned to the
agent, which embeds it in the now-complete InvocationBody and proceeds
with chain construction.

The approval signature is verified separately from the main aggregate.
This decoupling has two consequences. First, the security argument for
chain integrity does not depend on approval service trust:
compromise of an approval service does not enable chain forgery.
Second, approval verification adds a single pairing operation per
approval (typically zero or one per invocation), which we judge an
acceptable cost.

Policy can require approvals from multiple services (logical AND); the
required service identifiers and public keys are part of the policy
that the session token's \texttt{policy\_hash} commits to. An agent
cannot route approval to a different (perhaps more permissive) service
without invalidating the policy hash binding, defending against the
approval routing attack (T11).

\subsection{Verification}
\label{sec:protocol-verification}

The verifier algorithm proceeds in phases, ordered to fail fast on
cheap checks before performing expensive cryptographic operations.

\paragraph{Structural and temporal checks.}
The verifier first decodes the received chain and confirms canonical
encoding (rejecting non-canonical inputs to prevent encoding-ambiguity
attacks). It then checks that the current time is within the
invocation's validity window, that expiry times are monotonically
non-increasing along the chain (each hop's expiry cannot exceed its
parent's), and that the chain has minimum length one.

\paragraph{Replay protection.}
The verifier maintains a bounded-time nonce cache indexed by
$(\text{invoker public key}, \text{nonce})$. Chains whose key is
already present in the cache are rejected as replays. Successful
verifications add the key to the cache with TTL equal to the chain's
remaining validity window plus a clock-skew tolerance. The cache size
is bounded by the product of maximum verification throughput and the
maximum validity window.

\paragraph{Public key chain consistency.}
The verifier checks that the public key fields across consecutive
tokens agree: the session token's agent key equals the first delegation
token's delegator key, each delegation's delegatee key equals the next
hop's delegator key, and the last delegation's delegatee key equals the
invoker's key.

\paragraph{Scope monotonicity and policy compliance.}
The verifier checks that the sub-scope at each hop is a subset of the
scope at the prior hop, using the subset relation defined in
\cref{sec:policy}. It also evaluates the invocation against the final
scope, rejecting if the invocation falls outside the permitted scope.
It checks that the session token's policy hash matches the policy the
verifier has loaded.

\paragraph{Approval verification.}
If policy requires approvals, the verifier confirms that all required
approval services have signed valid receipts binding to the correct
invocation digest, and rejects the chain otherwise.

\paragraph{Aggregate signature verification.}
Finally, the verifier reconstructs $m_0, m_1, \ldots, m_N$ from the
received bodies (using position-determined tags) and checks the
aggregate via a multi-pairing:
\begin{equation}
\pair{g_1}{\sigma_{\text{agg}}} \overset{?}{=}
    \prod_{k=0}^{N} \pair{\pk_k}{\HashG(m_k)}.
\end{equation}
If this check succeeds, all integrity, parameter-binding, and chain
provenance properties hold simultaneously.

\paragraph{Cost.}
The aggregate check has a single multi-pairing structure with $N+2$
internal pairings. In modern BLS libraries (e.g., \texttt{blst}), the
Miller loops of these pairings share computation, and only a single
final exponentiation is required. The empirical cost is well
approximated by $\alpha + \beta N$ where $\alpha$ is the dominant final
exponentiation cost and $\beta$ is the per-Miller-loop cost. For chain
lengths $N \le 10$ (which we expect to cover the overwhelming majority
of practical deployments), the variation with $N$ is small relative
to $\alpha$, and total verification latency is dominated by the constant
term. This is significantly more favorable than verification approaches
requiring network round-trips per hop (e.g., OAuth token
introspection), but does not amount to truly asymptotic $O(1)$
verification; the honest characterization is ``compact aggregate
signatures with multi-pairing efficiency.''

\subsection{Encoding}
\label{sec:protocol-encoding}

All signed structures use deterministic CBOR encoding per RFC 8949
§4.2.1, with the following additional rules. Map keys are integers
rather than text strings, eliminating UTF-8 normalization
considerations and reducing encoding size. Float and tagged values are
disallowed in protocol structures except where explicitly required.
Group elements are encoded using the standard ZCash-format compressed
serialization defined for BLS12-381, and decoders perform subgroup
membership checks on every group element to defend against
small-subgroup attacks.

Verifiers reject any structure whose received encoding differs from
the canonical re-encoding of the decoded content. This precludes a
class of attacks based on adversarially-crafted non-canonical encodings
that parse correctly but hash to different values than the
implementation expects.

% ============================================================
\section{Identity Infrastructure}
\label{sec:identity}

The cryptographic protocol of \cref{sec:protocol} assumes that
verifiers can resolve identity-to-public-key bindings for all
participants in a chain. This section describes the federated identity
infrastructure providing those bindings, including registration with
proof-of-possession, public-key resolution, key rotation, and
revocation.

\subsection{Federated Registry Model}
\label{sec:identity-federated}

DelegationChain does not assume a global root of trust. Each
participating organization operates its own \emph{registry}, which
issues identity certificates for agents, signing services, and
approval services within that organization. Each registry holds a
long-lived root keypair, with the root public key serving as the trust
anchor for all certificates the registry issues.

Two organizations establish trust through one of three mechanisms,
chosen per deployment. \emph{Direct peering} exchanges registry root
keys out-of-band as part of a business agreement; this is the simplest
model and adequate for small numbers of partners. \emph{Federated
transparency} publishes registry root keys and certificate issuances
to a shared append-only log, modeled on Certificate Transparency
\citep{rfc9162} but adapted to DelegationChain's certificate format;
verifiers trust the log's commitments and obtain inclusion proofs
binding claimed keys to log entries. \emph{Hierarchical attestation}
introduces an attestation authority signing statements binding
registry roots to organizational identifiers, analogous to a root
CA. The protocol permits all three; we recommend direct peering or
federated transparency over hierarchical attestation, the latter
inheriting the trust-concentration concerns of Web PKI.

These mechanisms are not mutually exclusive. A production deployment
may use direct peering with frequent partners and a federated log for
the long tail.

\subsection{Identifiers and Registration}
\label{sec:identity-registration}

Identifiers are structured as
\texttt{organization-id:kind:local-id}, where \texttt{kind} is one of
\texttt{agent}, \texttt{signer}, or \texttt{approver}. The
\texttt{kind} tag prevents identifier confusion across roles. The
\texttt{organization-id} matches the issuing registry's identifier.

Registration produces a certificate binding an identifier to a public
key, signed by the registry. Certificates contain version
information, the identifier, the public key, the kind, the registry's
identifier and public key, issuance and expiry timestamps, a
not-before timestamp, a serial number, and the registry's signature.
Default certificate lifetimes are short --- 24 hours to 7 days for
agents and signing services, with longer lifetimes permitted only in
combination with active revocation (\cref{sec:identity-revocation}).
Short lifetimes reduce reliance on revocation by ensuring that
compromised credentials naturally expire within a bounded window.

\subsection{Proof of Possession}
\label{sec:identity-pop}

BLS aggregate signatures are vulnerable to rogue-key attacks unless
registrants prove possession of the secret keys corresponding to
registered public keys \citep{boneh2018compact}. Without
proof-of-possession (PoP), an adversary could register a public key
$\pk^* = \pk' - \sum_i \pk_i$ for chosen targets $\pk_i$ and a key
$\pk'$ for which the adversary holds the secret; valid signatures
$\sigma_i$ from the targets could then be combined with an adversary
signature to forge an aggregate that includes $\pk^*$.

DelegationChain requires every registration to include a PoP signature
over a registry-issued challenge. The challenge includes the
registrant's claimed identifier and public key, the kind tag, the
registry's identifier, a registry-issued fresh nonce (single-use,
short expiry), and a timestamp. The registrant computes a BLS
signature over this challenge with the secret key corresponding to the
public key being registered. The registry verifies the PoP signature
before issuing the certificate. The registry-issued nonce prevents
replay of PoPs captured from other contexts.

Registries additionally perform subgroup membership checks on
registered public keys, rejecting points that do not lie in the
prime-order subgroup of $\G{1}$. This check is non-trivial on BLS12-381
and has been the source of real-world implementation vulnerabilities.
Implementations should use library-provided subgroup checks (e.g.,
\texttt{blst\_p1\_in\_g1}) rather than reimplementing them.

\subsection{Public Key Resolution}
\label{sec:identity-resolution}

Verifiers resolve identifiers to public keys by querying the issuing
registry's resolution endpoint. The response includes the certificate
and, if the registry is integrated with a transparency log, an
inclusion proof. The verifier validates the certificate signature
under the registry's root public key (cached from prior interaction or
obtained out-of-band) and the inclusion proof against the log's
current signed tree head.

Verifiers cache resolved certificates to avoid registry queries on
every verification. Cache entries are valid until the earlier of the
certificate's expiry and a configurable cache TTL (typically one
hour). Cache invalidation is event-driven for revocation: observation
of a revocation assertion immediately evicts the corresponding cached
entry.

\subsection{Key Rotation}
\label{sec:identity-rotation}

The protocol supports three rotation modes. \emph{Scheduled rotation}
issues a new certificate with overlapping validity windows: during the
overlap, both old and new keys are valid, ensuring that chains in
flight at the moment of rotation continue to verify. Default rotation
cadence is every 30 days for agents and every 90 days for signing
services and approval services. \emph{Emergency rotation} truncates
the validity of the prior certificate via an immediately-issued
revocation assertion (\cref{sec:identity-revocation}) and issues a new
certificate without overlap. \emph{Registry root rotation} is the most
sensitive operation: the current root signs a rotation attestation
binding the new root to the old, which is then distributed via the
same channel used for the original root. Compromised-root recovery
requires out-of-band re-bootstrap of trust and cannot be handled by
in-band rotation.

When a hop's key rotates after that hop has signed but before the
chain reaches the verifier, the verifier uses the public key that was
valid at signing time, retrieved via a temporal resolution query that
returns the certificate valid at a specified timestamp.

\subsection{Revocation}
\label{sec:identity-revocation}

DelegationChain prefers short-lived certificates over revocation
infrastructure: with default certificate lifetimes of 24 hours to 7
days, the maximum window during which a compromised key can be misused
is bounded by certificate expiry. Explicit revocation is provided as a
defense-in-depth mechanism for deployments choosing longer-lived
certificates.

A \emph{revocation assertion} is a signed statement by a registry
declaring that a specific certificate is revoked, identified by serial
number. Verifiers learn of assertions through polling (default
interval 60 seconds), push notification, or transparency log
monitoring; the choice is a deployment decision. The end-to-end window
between key compromise and unforgeability is bounded by the sum of
detection latency, reporting latency, assertion issuance,
propagation, and residual cache time, of which the propagation and
cache terms are bounded by configuration (typically below 120
seconds).

Revocation applies relative to \emph{signing time}, not verification
time. A certificate revoked at time $t_r$ remains valid for chains
whose tokens were signed before $t_r$. Without this rule, every
revocation would retroactively invalidate all chains involving the
revoked identity, which would be operationally untenable.

% ============================================================
\section{Policy Expression Language}
\label{sec:policy}

This section introduces the policy expression language used to describe
agent authorization scopes throughout DelegationChain. Policies are
referenced by hash in session tokens (the \texttt{policy\_hash} field of
\cref{sec:protocol-digest}) and govern what each agent is permitted to
do, including what scope a delegating agent may transmit to a delegatee
and which invocations a final agent may make.

The design of the policy language is constrained by three requirements
arising from its role in the protocol. First, the verifier must be able
to decide, given two scope expressions $S_1$ and $S_2$ purportedly in a
delegation relationship, whether $S_2 \subseteq S_1$ --- the scope
monotonicity check of \cref{sec:protocol-verification}. Second,
evaluation of an invocation against a scope must terminate in time
bounded by the size of the scope expression and the invocation, with
no possibility of unbounded computation. Third, scope expressions must
be free of side effects: evaluation must not perform network calls,
file I/O, or other operations that could be exploited for information
exfiltration or denial of service.

These requirements rule out general-purpose policy languages such as
Rego or Cedar, which are powerful but operationally expensive to reason
about and historically have included features (e.g., HTTP fetches in
some versions of Rego) that violate the side-effect-free requirement.
DelegationChain's policy language is deliberately restricted to a
fragment that supports realistic enterprise authorization conditions
while admitting straightforward decidability and termination analyses.

\subsection{Grammar}
\label{sec:policy-grammar}

The grammar of scope expressions, in Backus-Naur form:

\begin{lstlisting}[basicstyle=\ttfamily\footnotesize]
scope    ::= rule | "allow" "all" | "deny" "all"
rule     ::= "allow" tool action params? where? approval?
tool     ::= "tool" "=" identifier
action   ::= "action" "=" identifier
params   ::= "params" "{" binding ("," binding)* "}"
binding  ::= path ":" type
where    ::= "where" expr
expr     ::= expr "and" expr
           | expr "or"  expr
           | "not" expr
           | "(" expr ")"
           | atom
atom     ::= path numop value
           | path strop string-value
           | path "in" list-value
           | path "under" path-value
type     ::= "int" | "string" | "bool" | "path"
numop    ::= "<" | "<=" | "==" | ">=" | ">"
strop    ::= "==" | "starts_with" | "ends_with" | "contains"
path     ::= identifier ("." identifier)*
value    ::= integer | string | boolean
list-value ::= "[" value ("," value)* "]"
\end{lstlisting}

A scope is either a special form (\texttt{allow all}, \texttt{deny all})
or a non-empty disjunction of \emph{rules}. Each rule specifies a tool,
an action, an optional declaration of expected parameter types
(\texttt{params}), an optional condition on parameter values
(\texttt{where}), and an optional approval requirement
(\texttt{approval}). The semantics of a scope is the union of the
invocations permitted by its rules; \texttt{deny all} permits nothing,
and \texttt{allow all} is a reserved form used for top-level trusted
contexts (e.g., the issuer's session-issuance authority).

The condition language inside \texttt{where} clauses is propositional:
boolean combinations of atomic predicates over parameter values. Each
atomic predicate is one of four forms --- numeric comparison, string
operator, list membership (\texttt{in}), or path-prefix containment
(\texttt{under}). No predicate involves function calls, network
operations, or unbounded iteration.

\subsection{Example}
\label{sec:policy-example}

A representative policy fragment governing a payments tool:

\begin{lstlisting}[basicstyle=\ttfamily\footnotesize]
allow tool=payments action=transfer
  params { amount: int,
           to_account: string }
  where amount <= 1000
allow tool=payments action=transfer
  params { amount: int,
           to_account: string }
  where amount <= 100000
        and to_account in
            ["acct_vendor_a", "acct_vendor_b"]
  approval requires approval_service_finance
allow tool=files action=read
  params { path: path }
  where path under "/home/agent/workspace"
\end{lstlisting}

This policy permits unconditional transfers up to 1{,}000 units;
transfers up to 100{,}000 units to two approved vendors with finance
approval; and read access to files anywhere under a specific workspace
directory. Any invocation falling outside these three rules is denied
by the implicit closing rule \texttt{deny all}.

\subsection{Evaluation Semantics}
\label{sec:policy-eval}

Given a scope $S$ and an invocation
$I = (\text{tool}, \text{action}, \text{params})$, the evaluator returns
\textsc{allow}, \textsc{allow-with-approval}(\textit{services}), or
\textsc{deny}. The procedure is:

\begin{enumerate}
    \item For each rule $r$ in $S$ in declaration order:
    \begin{enumerate}
        \item If $r$'s tool and action do not match $I$, continue.
        \item Type-check $I.\text{params}$ against $r$'s
              \texttt{params} declaration. If any parameter has the
              wrong type or any declared parameter is missing, continue.
        \item Resolve each path mentioned in $r$'s \texttt{where}
              clause against $I.\text{params}$. If any path resolves
              to absent, continue (treating the rule as inapplicable).
        \item Evaluate $r$'s \texttt{where} clause. If false, continue.
        \item If $r$ has no approval requirement, return \textsc{allow}.
              Otherwise return \textsc{allow-with-approval}($r$'s
              required approval services).
    \end{enumerate}
    \item If no rule matched, return \textsc{deny}.
\end{enumerate}

The evaluator is deterministic, has no side effects, and contains no
loops over data of unbounded size: the outer loop iterates over rules
(a finite syntactic component of $S$), and \texttt{where} clause
evaluation is a finite recursive descent over the expression tree.

\begin{proposition}[Termination]
\label{prop:termination}
Evaluation of an invocation against a scope $S$ terminates in time
$O(|S| \cdot |\text{params}|)$, where $|S|$ is the syntactic size of the
scope expression and $|\text{params}|$ is the total size of the
invocation's parameter encoding.
\end{proposition}

\begin{proof}[Proof sketch]
The outer loop iterates over rules of $S$, of which there are at most
$|S|$. For each rule, parameter resolution traverses the parameter map,
which is bounded by $|\text{params}|$; \texttt{where} expression
evaluation is a recursive descent whose depth is bounded by the
syntactic depth of the expression, itself bounded by $|S|$. The total
work is bounded by the product, yielding the claimed bound.
\end{proof}

\subsection{Subset Relation}
\label{sec:policy-subset}

The verifier of \cref{sec:protocol-verification} must check that the
sub-scope at each delegation hop is a subset of its parent's scope.
This is a syntactic check over scope expressions, not a runtime check
over invocations.

\begin{definition}[Subset relation]
\label{def:subset}
Given scope expressions $S_1, S_2$, we say $S_2 \subseteq S_1$ iff every
invocation $I$ that the evaluator accepts under $S_2$ is also accepted
under $S_1$, with at least as restrictive an approval requirement.
\end{definition}

The subset relation is decided by the following algorithm:

\begin{enumerate}
    \item If $S_1$ is \texttt{allow all}, return true.
    \item If $S_2$ is \texttt{deny all} or has no rules, return true.
    \item For each rule $r_2$ in $S_2$, attempt to find a rule $r_1$
          in $S_1$ such that $r_2 \subseteq r_1$ as rules. If no such
          $r_1$ exists for some $r_2$, return false.
    \item Otherwise, return true.
\end{enumerate}

Rule subsumption ($r_2 \subseteq r_1$) requires: the tools and actions
match; $r_1$'s parameter type declaration is compatible with $r_2$'s;
$r_1$'s approval requirement is no stronger than $r_2$'s (so that any
approvals satisfying $r_2$ also satisfy $r_1$); and $r_2$'s
\texttt{where} clause implies $r_1$'s.

The remaining sub-problem is deciding implication between propositional
formulas over the four atomic predicate forms. The atoms have a
restricted structure: numeric comparisons are over a fixed integer
variable and constants, string operators are over fixed string
variables and constant patterns, list membership is over fixed lists,
and path-prefix containment is over fixed prefixes. Implication
between such atoms is decidable in polynomial time by analyzing the
range or pattern relations on the underlying values. Implication
between propositional combinations is decidable by reduction to SAT
over the atomic predicates; since the predicates over distinct paths
are independent and predicates over the same path admit closed-form
implication checks, in practice the decision procedure runs in time
polynomial in the size of the scope expressions for the policies we
expect in real deployments.

\begin{proposition}[Decidability of subset]
\label{prop:decidability}
The relation $S_2 \subseteq S_1$ is decidable, with the procedure
running in time polynomial in $|S_1| + |S_2|$ for policies satisfying
mild structural conditions (no predicates that simultaneously constrain
the same path with both numeric and string operators in conflicting
ways).
\end{proposition}

\begin{proof}[Proof sketch]
The decision procedure described above performs a finite enumeration
over rule pairs, with each rule-pair check reducing to propositional
implication over decidable atomic predicates. Total work is bounded
polynomially in the input size.
\end{proof}

\subsection{Limitations of the Policy Language}
\label{sec:policy-limitations}

The policy language is deliberately bounded. It cannot express
semantic conditions of the form ``is this email message
professional?'' or ``is this code review constructive?'' that would
require external evaluation, typically by another LLM. Such conditions
are non-deterministic and admit no obvious decidability or termination
guarantee, and we judge them inappropriate for inclusion in
authorization policies that must be cryptographically committed to via
\texttt{policy\_hash}.

The language also cannot express conditions referring to state outside
the invocation itself: account balances, current time relative to
business hours, counts of prior invocations within a window. Such
conditions either require stateful evaluation (which would violate the
side-effect-free requirement) or external queries (which would
introduce latency and additional trust dependencies). Where stateful
authorization is required, we anticipate it being handled at the tool
implementation layer, after the protocol's stateless authorization
check has succeeded.

These limitations are deliberate. A more expressive language would
provide more flexibility but at the cost of guarantees we consider
essential. We discuss possible extensions to semantic conditions via
verifiable LLM attestation in \cref{sec:discussion}.

% ============================================================
\section{Security Analysis}
\label{sec:security}

This section presents the security argument for DelegationChain. We
state the protocol's security properties as theorems, reduce each to
standard cryptographic assumptions or to the properties of the policy
language, and sketch proofs. The arguments here are informal in the
sense that they have not been mechanized in a proof assistant or
verified using symbolic protocol analyzers such as Tamarin or
ProVerif. We discuss the planned formalization in
\cref{sec:limitations}.

\subsection{Assumptions}
\label{sec:security-assumptions}

The security arguments below rest on the following assumptions:

\begin{description}
    \item[\textbf{BLS unforgeability.}] The BLS signature scheme on
    BLS12-381 \citep{boneh2001bls,ietf-bls-signature} is
    existentially unforgeable under chosen-message attack, under the
    co-computational Diffie-Hellman assumption on the curve and in the
    random oracle model for $\HashG$.
    \item[\textbf{BLS aggregate security.}] Given proof-of-possession
    on all registered public keys, aggregate signatures are
    existentially unforgeable: no PPT adversary can produce a tuple
    $(\pk_1, \ldots, \pk_n, m_1, \ldots, m_n, \sigma_{\text{agg}})$
    such that the multi-pairing check accepts, unless the adversary
    has obtained valid individual signatures on each $m_i$ from the
    holder of $\sk_i$ \citep{boneh2018compact}.
    \item[\textbf{Collision resistance of SHA-256.}] SHA-256 is
    collision-resistant in the standard sense.
    \item[\textbf{Honest signing key custody.}] Signing keys remain in
    their designated hardware security modules and are not exfiltrated.
    Compromise of an HSM is explicitly out of scope per
    \cref{sec:threat-outside}.
    \item[\textbf{Honest registry root.}] Registry root signing keys
    are honestly managed. Compromise of a root is addressed via
    rotation and detection (\cref{sec:identity-rotation}); the
    formal arguments below assume an honest root, with the implication
    that root compromise reduces security to a recovery problem rather
    than a forgery problem.
\end{description}

\subsection{Theorems}
\label{sec:security-theorems}

\begin{theorem}[Token unforgeability]
\label{thm:unforgeability}
Under BLS unforgeability, no PPT adversary can produce a session,
delegation, or invocation token that verifies under a public key
$\pk$ for which it does not hold the corresponding $\sk$, except with
negligible probability.
\end{theorem}

\begin{proof}[Proof sketch]
A valid token consists of a body and a BLS signature on the body's
domain-tagged hash. Forging the signature without $\sk$ contradicts
BLS unforgeability under chosen-message attack: any adversary that
forges a token can be used to construct an adversary against BLS
unforgeability by treating the domain-tagged hash as the message and
extracting the signature.
\end{proof}

\begin{theorem}[Parameter binding]
\label{thm:parameter-binding}
A valid invocation token for parameters $P$ cannot be repurposed as
a valid invocation token for parameters $P' \neq P$, except with
negligible probability.
\end{theorem}

\begin{proof}[Proof sketch]
The signed message $m_N$ contains $\Canon(\text{InvocationBody})$,
which includes $\Canon(P)$. Under collision resistance of SHA-256 and
deterministic encoding, $\Canon(P) \neq \Canon(P')$ implies $m_N \neq
m_N'$. The signature $\sigma_N$ valid on $m_N$ is invalid on $m_N'$
by BLS unforgeability. The redundant \texttt{params\_hash} field
provides defense-in-depth against canonical encoding bugs but is not
required for the theorem.
\end{proof}

\begin{theorem}[Chain integrity]
\label{thm:chain-integrity}
The aggregate signature over a delegation chain covers the entire
chain content: any modification, reordering, truncation, or extension
of the chain invalidates the aggregate, except with negligible
probability.
\end{theorem}

\begin{proof}[Proof sketch]
Modification of any body changes its canonical encoding, changing the
corresponding $m_k$ (since $m_k$ hashes the body). Reordering changes
the sequence of $m_k$ values (since $m_k$ commits to $m_{k-1}$).
Truncation removes a contribution from the aggregate, which then fails
the multi-pairing check against the original set of public keys.
Extension requires producing a new $\sigma$ for the extension hop;
under BLS aggregate security and proof-of-possession, this requires
the secret key of the extending party. The position-determined tag
check (\cref{sec:protocol-digest}) rejects chains where an invocation
token is repurposed as a delegation token (or vice versa) for chain
extension.
\end{proof}

\begin{theorem}[Scope monotonicity]
\label{thm:scope-monotonicity}
In any chain accepted by the verifier, the scope at hop $k+1$ is a
subset of the scope at hop $k$, in the sense of
\cref{def:subset}.
\end{theorem}

\begin{proof}[Proof sketch]
The verifier checks $\text{sub\_scope}_{k+1} \subseteq
\text{scope}_k$ explicitly (\cref{sec:protocol-verification}) using
the decision procedure of \cref{sec:policy-subset}. Acceptance of the
chain implies this check passed. The decidability of the subset
relation (\cref{prop:decidability}) makes the check well-defined.
The integrity of the body content from which sub-scopes are read is
guaranteed by \cref{thm:chain-integrity}.
\end{proof}

\begin{theorem}[Replay resistance]
\label{thm:replay}
A verifier that maintains a correctly-sized nonce cache and rejects
chains whose nonce-and-invoker-key pair is already in the cache will
reject all replays within the validity window, and will reject all
chains presented after the validity window has expired.
\end{theorem}

\begin{proof}[Proof sketch]
Within the validity window, the nonce cache contains the
$(\pk_N, \text{nonce})$ pair of every previously accepted chain. A
replay presents the same pair and is rejected. Outside the validity
window, the verifier rejects on the expiry check before performing
signature verification. The cache TTL is set to the validity window
plus clock-skew tolerance, ensuring no premature eviction.
\end{proof}

\begin{theorem}[Approval binding]
\label{thm:approval-binding}
An approval receipt for invocation $I$ cannot be reused for any
invocation $I' \neq I$, and an invocation requiring approval from
service $s$ cannot be accepted with approval from service $s' \neq s$,
except with negligible probability.
\end{theorem}

\begin{proof}[Proof sketch]
The approval body signed by the approval service includes
\texttt{invocation\_digest}, a hash of the invocation body without the
approval reference itself. Modification of the invocation changes
this digest, invalidating the receipt under BLS unforgeability. The
required approval service identifiers are committed to by the policy
hash referenced in the session token, which the chain cannot
forge; routing to a different service produces a chain whose policy
hash does not match the verifier's loaded policy, and is rejected.
\end{proof}

\subsection{Threats Addressed}
\label{sec:security-threats}

The seventeen threats of \cref{sec:threat-model} are addressed by the
above theorems and the identity infrastructure of
\cref{sec:identity}, as follows. Token forgery (T1) is addressed by
\cref{thm:unforgeability}. Parameter substitution (T2) is addressed by
\cref{thm:parameter-binding}. Scope escalation (T3) and chain forgery
(T4) are addressed by Theorems \ref{thm:chain-integrity} and
\ref{thm:scope-monotonicity}. Cross-session replay (T6) and intra-window
replay (T7) are addressed by \cref{thm:replay}. Cross-context token
reuse (T9) is addressed by domain separation tags
(\cref{sec:protocol-digest}). Approval-related threats (T10--T13) are
addressed by \cref{thm:approval-binding} and the policy-hash binding.
Identity threats (T14--T17) are addressed by proof-of-possession
(\cref{sec:identity-pop}), short-lived certificates and revocation
(\cref{sec:identity-revocation}), and optional transparency log
integration.

Prompt injection within policy (T5) and confused-deputy delegation (T8)
require separate discussion. \emph{Within-policy prompt injection} ---
an adversary inducing an agent to take an action that is within its
permitted policy but contrary to the principal's intent --- is not
prevented by the protocol. The protocol's contribution to defending
against this threat is to \emph{bound its blast radius}: the adversary
can cause only invocations that policy already permits and, where
policy requires human approval, will surface the action to an approver
who has an opportunity to refuse. Tightly-scoped policies and
appropriately-required approvals reduce the blast radius; defending
against prompt injection itself remains an active research area
outside the protocol's purview. \emph{Confused-deputy delegation},
where an agent is induced to delegate authority to an adversary, is
similarly bounded rather than prevented: the agent cannot delegate
more authority than it itself holds (\cref{thm:scope-monotonicity}),
limiting damage to the agent's own permitted scope.

\subsection{Limitations of the Analysis}
\label{sec:security-limitations-analysis}

We emphasize that the analysis above consists of proof sketches under
specified assumptions, not machine-checked or peer-reviewed proofs.
Several aspects warrant scrutiny in any subsequent formalization.
\Cref{thm:chain-integrity} rests on the assumption that
proof-of-possession is correctly enforced at registration; an
implementation that fails to verify PoP would invalidate the chain
integrity guarantee against rogue-key attacks. \Cref{thm:replay}
depends on operational properties of the nonce cache (correct sizing,
crash recovery, distributed consistency across verifier replicas)
that are not enforced by the protocol but must be discharged by the
verifier implementation. The subset relation algorithm of
\cref{sec:policy-subset} runs in polynomial time under mild
structural conditions on policy expressions; adversarially-crafted
policies could exercise corner cases requiring closer analysis.
Symbolic verification of the protocol in Tamarin or ProVerif would
confirm or refute these arguments and is planned future work.

% ============================================================
\section{Discussion}
\label{sec:discussion}

This section discusses related work, limitations of the protocol and
its analysis, and directions for future work.

\subsection{Related Work}
\label{sec:related-work}

DelegationChain draws on three distinct research traditions:
capability-based authorization, pairing-based cryptographic protocols,
and agent-system security. We position our contribution relative to
each.

\paragraph{Capability-based authorization.}
The conceptual framework of attenuated capabilities, in which authority
is represented by unforgeable tokens that can be progressively
restricted as they are delegated, dates to early work on capability
operating systems and was formalized in the modern context by
\citet{miller2006robust}. The most directly relevant prior work is
Macaroons \citep{birgisson2014macaroons}, which applied attenuated
capabilities to cloud authorization using HMAC chains. A Macaroon
carries ``caveats,'' predicates that restrict its use; each delegation
appends additional caveats, and verification recomputes the HMAC chain
using a shared secret known to the verifier. Macaroons demonstrated
that attenuated delegation could be made cryptographic and efficient.
DelegationChain extends this model in three principal directions: from
shared-secret HMAC to public-key BLS signatures (enabling
cross-organizational verification without shared secrets); from
per-hop linear signature growth to compact aggregate signatures via
BLS aggregation; and from informal caveat languages to a typed
declarative policy language with decidability guarantees. Where
Macaroons assumes a single trust domain sharing HMAC keys, our work
targets the cross-organizational setting where mutual distrust is the
norm.

UCAN \citep{ucan-spec} also extends capability delegation to a
public-key setting, originally for Web3 contexts. UCAN uses ECDSA or
Ed25519 signatures and a JWT-like token format, with delegation
attenuation expressed through scoped capability tokens. Our work
differs from UCAN in several respects: we use BLS signatures
specifically for their aggregation property, which UCAN does not
provide; we specify a typed policy language with formal decidability
of subset, where UCAN's capability expressions are less rigorously
defined; we target MCP rather than Web3 transport; and we include a
federated identity infrastructure with proof-of-possession and
revocation semantics adapted to the short-lived-credential model.

\paragraph{Cryptographic infrastructure.}
The BLS signature scheme is due to \citet{boneh2001bls}. Aggregate
signature security under proof-of-possession is analyzed by
\citet{boneh2018compact}, whose work motivates our PoP requirement
during registration. We use the IETF-standardized BLS12-381
curve and signing variant \citep{ietf-bls-signature} and the
hash-to-curve construction of \citep{rfc9380}. Our chain digest
construction is, to our knowledge, novel: while recursive hash chains
are a standard cryptographic primitive, their composition with BLS
aggregation to provide simultaneous chain integrity and bandwidth
compactness for delegation chains has not previously been specified.
Certificate Transparency \citep{rfc9162} provided the model for our
optional transparency log integration.

\paragraph{MCP and agent-system security.}
The security gap in cross-organizational MCP deployments has been
empirically characterized by recent measurement work
\citep{huang2026inch}, finding that a substantial fraction of
deployed MCP servers lack adequate caller identity verification.
\citet{maloyan2026breaking} identify additional protocol-level
vulnerabilities in MCP. These works establish the problem but do not
propose cryptographic solutions. Single-organization runtime
enforcement approaches \citep{buhler2025agentbound} provide
defense-in-depth within an organization's perimeter but do not
provide cross-organizational cryptographic proofs and are not
designed to extend across organizational boundaries.

The closest contemporary work is the Agent Identity Protocol of
\citet{prakash2026aip}, which proposes invocation-bound capability
tokens for cross-protocol agent delegation. AIP and DelegationChain
share the high-level goal of providing verifiable invocation
authorization in agent settings. The principal differences are:
DelegationChain provides BLS-based aggregate signatures over
multi-hop chains, where AIP relies on individual signatures per hop;
DelegationChain specifies a formal policy expression language with
decidability proofs, where AIP's policy semantics are less formally
treated; and DelegationChain includes a complete federated identity
infrastructure specification, where AIP focuses on the token
construction. The two designs are largely complementary, and either
could be adopted in deployments depending on the priority placed on
aggregate-signature compactness versus simplicity of individual
signature verification.

\paragraph{Workload and identity infrastructure.}
SPIFFE/SPIRE \citep{spiffe} provides workload identity through
short-lived X.509 certificates and JWT-format SVIDs. While SPIFFE
addresses workload identity within a trust domain, it does not
provide parameter-bound authorization, delegation chain integrity,
or aggregate signatures. OAuth 2.1 \citep{oauth21} addresses caller
identity and scope but requires synchronous token introspection on the
verification path and does not provide parameter binding. We view
DelegationChain not as a replacement for these systems but as
complementary infrastructure for the specific use case of verifiable
agent-to-tool invocation across organizations.

\subsection{Limitations}
\label{sec:limitations}

We discuss limitations of the work along several dimensions.

\paragraph{Analytical, not empirical.}
This paper presents a protocol design and preliminary security
analysis. It does not present empirical performance measurements from
a reference implementation. The verification cost analysis of
\cref{sec:protocol-verification} is based on counting cryptographic
operations and library-reported costs for those operations; precise
wall-clock latency in the deployment configurations of interest will
depend on implementation choices (BLS library, multi-pairing strategy,
deserialization performance) that this paper does not commit to.
A Rust reference implementation using the \texttt{blst} library is
under development; empirical performance data will appear in a
follow-up paper or extended version.

\paragraph{Security arguments are sketches, not proofs.}
\Cref{sec:security} presents proof sketches that reduce protocol
security to standard cryptographic assumptions. These sketches have
not been mechanized in a proof assistant nor have they been verified
by independent cryptographer review. Symbolic protocol verification
using Tamarin or ProVerif is planned and is expected to reveal at
minimum the kind of subtle interaction issues that hand-analysis
typically misses (e.g., key-compromise impersonation paths, unknown
key-share variants). We treat the current analysis as a starting
point for formal verification, not as a substitute for it.

\paragraph{Policy language expressiveness.}
The policy language of \cref{sec:policy} is deliberately bounded. It
cannot express semantic conditions requiring evaluation by an LLM or
external system (``is this email professional?''), stateful
conditions referring to data outside the invocation (``is this user
under their monthly limit?''), or conditions requiring synchronous
queries to external systems. These limitations are intentional: each
relaxation would compromise the decidability, termination, or
side-effect-freedom guarantees on which the security analysis depends.
Where richer policy expressions are required in practice, we
anticipate them being enforced at the tool implementation layer
\emph{after} the protocol's stateless check has succeeded. A
fully-integrated treatment of stateful policy may require extensions
to the protocol --- including verifiable LLM attestation for
semantic conditions --- which we leave as future work.

\paragraph{Threat model boundaries.}
Several classes of threats are explicitly out of scope
(\cref{sec:threat-outside}), including fully-compromised HSMs,
compromised signing service infrastructure, side-channel attacks,
and attacks on human approvers' devices. The protocol provides no
defense against these and relies on operational measures (HSM
hardening, key custody discipline, secure approver workflows) for
mitigation. Within-policy prompt injection is bounded but not
prevented (\cref{sec:security-threats}); the protocol's contribution
is to ensure that prompt-injected agents cannot exceed their permitted
policy, not to prevent prompt injection itself.

\paragraph{Post-quantum security.}
BLS12-381 is not post-quantum secure. A sufficiently capable quantum
adversary could forge BLS signatures. Migration to a post-quantum
signature scheme is a future protocol version; this likely involves
hybrid construction combining BLS with a post-quantum scheme during
a transition period, with the complexity tradeoffs that hybrid
constructions imply.

\paragraph{Bootstrapping and adoption.}
The federated registry model requires both communicating organizations
to operate compatible registries (or to participate in a shared
transparency log). Adoption of the protocol in cross-organizational
agent ecosystems thus faces a bootstrapping problem: early adopters
have limited peers to communicate with, and the value of adoption
grows with peer density. We do not address bootstrapping strategies
in this paper, but observe that the same dynamic affects all
cross-organizational identity protocols, including SPIFFE and
SAML federation. Successful adoption is likely to require either
integration into a widely-deployed platform (e.g., as a feature of
an MCP-compliant agent framework) or coordination through standards
bodies.

\subsection{Future Work}
\label{sec:future-work}

We identify several directions for future work, in approximate order
of priority.

\emph{Reference implementation and empirical evaluation.} A Rust
reference implementation using \texttt{blst} is the most immediate
next step. Empirical evaluation should measure end-to-end verification
latency across chain lengths, throughput under load, and storage
overhead for nonce caches and certificate caches. Comparison with
OAuth token introspection and SPIFFE SVID verification on
representative workloads will substantiate the analytical cost
arguments of \cref{sec:protocol-verification}.

\emph{Symbolic verification.} Modeling the protocol in Tamarin or
ProVerif and proving the security properties of \cref{sec:security}
will validate (or, more usefully, falsify) the hand-written proof
sketches. This is expected to surface protocol details requiring
revision before publication of a normative specification.

\emph{Differential privacy on audit trails.} The protocol generates
detailed audit records of every accepted invocation. For deployments
in regulated industries, the ability to produce compliance attestations
(``$X\%$ of invocations were authorized'') without exposing the
content of individual invocations would have practical value.
Standard differential privacy mechanisms apply; their integration into
the protocol's audit format is straightforward but was deferred from
this work for scope reasons.

\emph{Verifiable LLM attestation for semantic policy conditions.}
The most significant restriction of the policy language is its
inability to express semantic conditions. Recent work on verifiable
LLM inference (zkML and related approaches) may enable a future
extension in which semantic policy conditions are evaluated by an LLM
whose evaluation is itself cryptographically attested. This is
speculative; the practicality of zkML at the latencies required for
authorization protocols is not currently established.

\emph{Post-quantum migration.} As discussed above, planning for
migration to post-quantum signatures is appropriate even before the
threat is acute. Investigating hybrid BLS-plus-post-quantum
constructions adapted to the aggregation property is an open question.

\emph{Cross-registry naming and discovery.} The current spec does not
address how organizations choose and namespace their
\texttt{organization-id} values. In practice deployments will use
DNS-style hierarchical naming or other conventions, but standardizing
this is needed before adoption at scale.

% ============================================================
\section{Conclusion}
\label{sec:conclusion}

The growing deployment of LLM-based agents across organizational
boundaries creates an authorization problem that existing protocols do
not adequately address. Receiving organizations need cryptographic
assurance that an incoming tool invocation was authorized by a
specific agent operating within a specific scope on behalf of specific
principals, with specific parameter values --- and they need this
assurance without depending on the sending organization's
infrastructure being available or trustworthy at verification time.

We have presented DelegationChain, a protocol providing four
properties jointly: cross-organizational verification without a
shared trust anchor, parameter-bound authorization, aggregatable
proof chains, and offline verification. The protocol's cryptographic
core is a recursive chain digest construction over BLS12-381
aggregate signatures, with proof-of-possession-protected registration
defending against rogue-key attacks. The associated policy expression
language is decidable, terminating, and side-effect-free by
construction, supporting the verifier's scope monotonicity check
without introducing computational or semantic hazards. The federated
identity infrastructure permits interoperation between mutually
distrusting organizations through direct peering, federated
transparency logs, or hierarchical attestation, as deployment context
requires.

We have stated and sketched proofs of the protocol's core security
properties --- unforgeability, parameter binding, chain integrity,
scope monotonicity, replay resistance, and approval binding --- and
mapped these to the threat model. We have been explicit about the
threats the protocol does not address, particularly within-policy
prompt injection, which the protocol bounds rather than prevents.

The work has clear limitations: the analysis is hand-written rather
than mechanically verified, no empirical performance data is yet
available, and several threat classes remain out of scope. A
reference implementation in Rust, symbolic verification in Tamarin
or ProVerif, and empirical evaluation are the immediate next steps.
These additions would establish DelegationChain as a viable
foundation for the inter-agent authorization layer in production
deployments.

We view this paper as a starting point. The design space for
cross-organizational agent authorization is still being mapped, and
multiple competing proposals --- AIP, AGENTBOUND, MCP-Guard, and
others --- are appearing simultaneously with this work. We expect
useful integration and competition between these proposals over the
next several years as the inter-agent ecosystem matures. Our
contribution is one design point in this space, optimized for
combining attenuated capability semantics, public-key
cross-organizational verification, and aggregate-signature
efficiency. Whether this combination proves to be the right one for
the deployments that emerge is a question only the deployments will
answer.

% ============================================================
\section*{Acknowledgments}
\label{sec:acknowledgments}

The author thanks colleagues at EKTU for early
discussion of the protocol's design, and the broader MCP and agent
security research communities for the context in which this work
was undertaken. Detailed acknowledgments will be added in a
subsequent revision following peer review of the analysis.

\section*{Use of AI Tools}
\label{sec:ai-tools}

The author used a large language model (Anthropic Claude) extensively
during the preparation of this manuscript. The model was used to
draft prose, propose organizational structures for technical sections,
formalize the policy language grammar, and prepare initial proof
sketches for the security theorems. The model also served as a
critical reader, identifying issues such as overstated claims about
verification complexity that were subsequently corrected.

The protocol design, threat model, and all underlying research
decisions are the author's own. All technical claims have been
independently verified by the author against authoritative sources;
the model's outputs were not accepted without scrutiny. All cited
references were verified against original sources. The author takes
full responsibility for the content, accuracy, and conclusions of
this work.

\printbibliography

\end{document}